\chapter{Herleitung des Informationsflusses}
\label{anhang:fluss}

\section{Einleitung}

Die fundamentale Dynamik der IWT wird durch die diskrete Euler-Lagrange-Gleichung bestimmt. 
In der numerischen Implementierung wird der Fluss der Information jedoch häufig in der Form
\( J = K \cdot \Delta \rho \cdot \Delta Q \) dargestellt. 

Dieser Anhang zeigt die vollständige und konsistente Herleitung dieser Flussgleichung 
aus den Axiomen der IWT, insbesondere aus dem Variationsprinzip (Axiom 4) und der 
Kontinuitätsgleichung, die sich aus der Phasen-Amplitude-Darstellung von \( I \) ergibt.

Die Herleitung macht explizit, warum der Fluss nicht durch reine Diffusion (\( \nabla \rho \)) 
getrieben wird, sondern durch die Kopplung von Dichte und Bohm-Potential.

\section{Ausgangspunkt: Die diskrete Euler-Lagrange-Gleichung}

Aus Axiom 4 (Variationsprinzip) folgt die diskrete Euler-Lagrange-Gleichung:

\[
\frac{\partial \mathcal{L}_d}{\partial I_k^{(n)}} 
- \Delta_t^+ \left( \frac{\partial \mathcal{L}_d}{\partial (\Delta_t I_k^{(n)})} \right) 
- \Delta_s \left( \frac{\partial \mathcal{L}_d}{\partial (\Delta_s I_k^{(n)})} \right) = 0
\]
\label{eq:euler_lagrange}

Die Dynamik des Informationsfeldes \( I_k^{(n)} \) wird somit durch ein Wirkungsprinzip 
festgelegt. Die räumliche Kopplung ist im diskreten Laplace-Operator 
\( \Delta_s \) enthalten, der über die fraktale Kopplungsmatrix \( K_{kl} \) definiert ist.

\section{Die Kontinuitätsgleichung aus der Phasen-Amplitude-Darstellung}

Die fundamentale Größe der IWT ist das komplexe Feld:

\[
I_k^{(n)} = \sqrt{\rho_k^{(n)}} \cdot e^{i \phi_k^{(n)}}
\]

wobei \( \rho_k^{(n)} = |I_k^{(n)}|^2 \) die Informationsdichte (Wahrscheinlichkeitsdichte) 
und \( \phi_k^{(n)} = \arg(I_k^{(n)}) \) die Phase ist.

Die Separation der Euler-Lagrange-Gleichung in Real- und Imaginärteil (siehe Kapitel 3.4.3) 
ergibt zwei fundamentale Gleichungen. Der Imaginärteil liefert die diskrete Kontinuitätsgleichung:

\[
\frac{\rho_k^{(n+1)} - \rho_k^{(n)}}{T} 
+ \sum_{l \in \mathcal{N}(k)} K_{kl} \cdot \left( \rho_k^{(n)} \vec{v}_k^{(n)} - \rho_l^{(n)} \vec{v}_l^{(n)} \right) = 0
\]
\label{eq:kontinuitaet}

Dies ist die diskrete Form der Kontinuitätsgleichung:
\[
\frac{\partial \rho}{\partial t} + \nabla \cdot (\rho \vec{v}) = 0
\]

Die Geschwindigkeit \( \vec{v} \) ist aus dem Realteil der Euler-Lagrange-Gleichung als Gradient 
der Phase identifiziert:

\[
\vec{v}_k^{(n)} = \frac{\hbar}{m} \nabla \phi_k^{(n)}
\]
\label{eq:fuehrung}

\section{Die Kopplung von Phase und Bohm-Potential}

Der Realteil der Euler-Lagrange-Gleichung ergibt die Hamilton-Jacobi-Gleichung 
mit dem Bohm-Potential \( Q \):

\[
\frac{\partial S}{\partial t} + \frac{|\nabla S|^2}{2m} + V + Q = 0
\]
mit \( S = \hbar \phi \) und 
\[
Q = -\frac{\hbar^2}{2m} \frac{\nabla^2 \sqrt{\rho}}{\sqrt{\rho}}
\]

Im stationären Fall oder für langsame Veränderungen dominiert das Bohm-Potential 
die Dynamik der Phase. Insbesondere gilt für den Gradienten der Phase:

\[
\nabla \phi \propto \nabla Q
\]
\label{eq:phi_q_gradient}

Dies ist die zentrale Kopplung: Der Gradient des Bohm-Potentials \( Q \) bestimmt die Änderung der Phase.

\section{Der Fluss im diskreten Netzwerk}

Aus der Kontinuitätsgleichung (Gl. \ref{eq:kontinuitaet}) und der Führungsgleichung 
(Gl. \ref{eq:fuehrung}) folgt für den Fluss:

\[
\vec{J}_k^{(n)} = \rho_k^{(n)} \vec{v}_k^{(n)} 
= \rho_k^{(n)} \frac{\hbar}{m} \nabla \phi_k^{(n)}
\]

Im diskreten Fall wird der Gradient durch die Differenz zwischen benachbarten Knoten 
und die Kopplungsmatrix \( K_{kl} \) ersetzt:

\[
J_k^{(n)} = \sum_{l \in \mathcal{N}(k)} K_{kl} \cdot \left( \rho_k^{(n)} - \rho_l^{(n)} \right) \cdot \left( \phi_k^{(n)} - \phi_l^{(n)} \right)
\]
\label{eq:fluss_phi}

Dabei ist \( K_{kl} \) die fraktale Kopplungsmatrix aus Axiom 3:
\[
K_{kl} = \frac{1}{d_{kl}^{3-D}}
\]

\section{Substitution von \( \Delta \phi \) durch \( \Delta Q \)}

Die Beziehung \( \nabla \phi \propto \nabla Q \) (Gl. \ref{eq:phi_q_gradient}) 
erlaubt die Substitution der Phasendifferenz durch die Differenz des Bohm-Potentials:

\[
\phi_k - \phi_l \propto Q_k - Q_l
\]

Eingesetzt in Gl. \ref{eq:fluss_phi} ergibt dies:

\[
J_k^{(n)} = \sum_{l \in \mathcal{N}(k)} K_{kl} \cdot \left( \rho_k^{(n)} - \rho_l^{(n)} \right) \cdot \left( Q_k^{(n)} - Q_l^{(n)} \right)
\]
\label{eq:fluss_q}

Dies ist die gesuchte Flussgleichung, die in der numerischen Implementierung verwendet wird.

\section{Die physikalische Interpretation}

Gl. \ref{eq:fluss_q} hat eine tiefgehende physikalische Bedeutung:

\begin{enumerate}
    \item \textbf{Keine reine Diffusion:} Der Fluss wird nicht allein durch den Dichtegradienten 
    \( \Delta \rho \) getrieben. Würde nur Diffusion wirken, wäre der Fluss \( J \propto \Delta \rho \), 
    was zu einem vollständigen Zerfall aller Strukturen führen würde. Die Strukturbildung 
    wäre unmöglich.

    \item \textbf{Kopplung von Dichte und Potential:} Der Fluss ist das Produkt aus dem 
    Dichteunterschied und dem Potentialunterschied. Dies reflektiert die Selbstorganisation: 
    Information fließt dorthin, wo sowohl ein Dichte- als auch ein Potentialgefälle existiert.

    \item \textbf{Stabilität:} An Stellen, an denen die Dichte extremal ist (\( \Delta \rho = 0 \)), 
    verschwindet der Fluss. Ebenso verschwindet er, wo das Potential extremal ist 
    (\( \Delta Q = 0 \)). Dadurch entstehen stabile, lokalisierte Strukturen.

    \item \textbf{Bindende Wirkung von \( Q \):} Da \( Q \) im Zentrum von Strukturen negativ 
    ist (minimum) und nach außen ansteigt, ist der Fluss an den Rändern maximal und zeigt 
    nach innen. Im Zentrum verschwindet er jedoch, sodass die Struktur kollapsfrei bleibt.
\end{enumerate}

\section{Die Implementierung im Code}

Die Gleichung \ref{eq:fluss_q} wird im OpenCL-Kernel \texttt{iwt\_flux} wie folgt umgesetzt:

\begin{verbatim}
sum += Kij * (rho_i - rho_j) * (Q_i - Q_j);
\end{verbatim}

Dabei ist:
\begin{itemize}
    \item \texttt{Kij} – die fraktale Kopplung zwischen Knoten \( i \) und \( j \)
    \item \texttt{rho\_i} – \( \rho_i = |I_i|^2 \)
    \item \texttt{Q\_i} – das Bohm-Potential am Knoten \( i \)
\end{itemize}

Diese Implementierung ist keine willkürliche Wahl, sondern die exakte diskretisierte Form 
der aus der Theorie abgeleiteten Flussgleichung.

\section{Zusammenfassung}

\begin{enumerate}
    \item Der Fluss \( J = K \cdot \Delta \rho \cdot \Delta Q \) folgt aus der 
    diskreten Euler-Lagrange-Gleichung (Axiom 4) und der Kontinuitätsgleichung.
    \item Die Kopplung zwischen Phase \( \phi \) und Bohm-Potential \( Q \) 
    wird durch die Hamilton-Jacobi-Gleichung hergestellt.
    \item Die fraktale Kopplungsmatrix \( K_{kl} = 1/d_{kl}^{3-D} \) ist die 
    emergente Metrik des Raumes (Axiom 3).
    \item Die Implementierung im Code ist keine Näherung, sondern die exakte 
    diskrete Form der fundamentalen Dynamik der IWT.
\end{enumerate}

Die Herleitung zeigt, dass die beobachtete Stabilität der Strukturen und das 
Ausbleiben eines gravitativen Kollapses keine numerischen Artefakte sind, 
sondern eine direkte Konsequenz der mathematischen Struktur der IWT.